\scanpage{9}
\uline{Lemme 6} \quad Soient $S$ schéma \deleted{\unsure{de}} \inserted{$\mathrm{loc}^{\unsure{t|.}}$} noeth, $F$ \deleted{\illegible{1}} coh. sur $S$.

a) Soit $J$ le Nilradical de $\mathcal O_S$, \inserted{$S_0=S_{\mathrm{réd}}$.} Il existe un ouvert dense $U$ de $S$ tel que \deleted{\illegible{3} \unsure{et|(2)} $\unsure{\operatorname{gr}|\operatorname{Supp}}_{\!J}(F)$} $\operatorname{gr}_J(\mathcal O_S)$ \inserted{et $\operatorname{gr}_J(F)$} loc. libres sur $U_0$

b) Soit $U$ comme dessus. Alors pour tout $s\in U$, et \deleted{toute} suite $\mathcal O_{S,s}$-régulière $f_1\ldots f_n$ \deleted{de} $\unsure{\mathfrak m_s|M_s}$, \deleted{posant $M=\mathcal O_{S,s}/(f_1,\ldots,f_n)$,} on a
\[
\operatorname{Tor}^{\mathcal O_{S,s}}_i\bigl(\deleted{\unsure{\mathcal O_{S,s}}}F_s,\mathcal O_{S,s}/(f_1,\ldots,f_n)\bigr)=0\qquad\text{pour }i>0.
\]

\uline{Dém} L'assertion a) est triviale. Pour b), on est ramené au cas où $S=U$, \deleted{et $F$ provient d'un faisceau loc libre sur $S_0$.} \inserted{et où $S=\operatorname{Spec}\mathcal O_{S,s}$,} Pour calculer les $\operatorname{Tor}_i$ on utilise le complexe de chaines de Koszul $K_{\bullet}(\underrightarrow f)$, qui est une résolution de $A/(f_1\ldots,f_n)$ $(A=\mathcal O_{S,s})$. On est ramené à prouver que les images $f_{10}\ldots f_{n0}$ des $f_i$ de $A_0$ forment \deleted{\unsure{encore|aussi}} \inserted{aussi} une suite $M$-régulière. En fait, on prouve en $=$ temps que, posant $N=M/(f_1\ldots f_n)M$, on a $\operatorname{gr}_J(N)$ \deleted{loc. libre} libre. C'est bien standard si $n=1$, et on continue par récurrence.

\uline{Corollaire} 7 \quad $S$ schéma noethérien, $P$ et $\Omega$ faisceaux cohérents. Alors $\exists$ ouvert dense \inserted{$U$} de $S$ tel que pour tt $s\in U$ et tt suite $\underrightarrow f$ $\mathcal O_{S,s}$-régulière $f_1\ldots f_n$,

% archive sheet number: 8
